\documentclass[11pt]{article}

% ---------- Packages ----------
\usepackage[utf8]{inputenc}
\usepackage[margin=1in]{geometry}
\usepackage{amsmath, amssymb, amsthm, mathtools}
\usepackage{microtype}
\usepackage{enumitem}
\usepackage{booktabs}
\usepackage{graphicx}
\usepackage{xcolor}
\usepackage{hyperref}
\hypersetup{
  colorlinks=true,
  linkcolor=blue!60!black,
  citecolor=blue!60!black,
  urlcolor=blue!60!black
}

% ---------- Helpers ----------
\newcommand{\refnum}[1]{\textsuperscript{[#1]}}
\newcommand{\R}{\mathbb{R}}

% ---------- Title ----------
\title{\vspace{-1em}\textbf{Conceptual Bridges in Machine Learning: From Foundations to Frontiers}}
\author{}
\date{}

\begin{document}
\maketitle

\section*{Introduction}
The modern landscape of machine learning and statistics presents a vast and often bewildering array of algorithms, models, and architectures. From the transparent simplicity of linear regression to the labyrinthine complexity of deep generative models, the field can appear as a disconnected collection of specialized tools. However, a deeper examination reveals a structured and interconnected web of ideas, where advanced concepts often emerge as principled extensions and generalizations of foundational ones. Understanding these evolutionary pathways---the ``conceptual bridges'' that link simple models to their more complex counterparts---is fundamental to achieving a nuanced mastery of the discipline. This report provides an exhaustive exploration of these connections, illustrating how a solid grasp of the basics serves as the most reliable guide to the frontiers of research.

The evolution from simple to complex models is not random but is driven by a consistent set of powerful themes. By repeatedly applying a few core principles, the field has systematically overcome the limitations of earlier methods to build progressively more powerful and flexible tools. The primary recurring themes of this evolution, which will be explored throughout this report, include:
\begin{itemize}[leftmargin=1.2em]
  \item \textbf{Introducing Non-Linearity:} The transition from models that assume linear relationships to those capable of capturing the complex, curved, and intricate patterns inherent in real-world data. This is often achieved by composing simple linear operations with non-linear activation functions.
  \item \textbf{Embracing Probability:} The shift from deterministic models that provide a single point estimate to probabilistic models that represent uncertainty through distributions. This allows for a more complete and robust understanding of data, enabling tasks like generative modeling and uncertainty quantification.
  \item \textbf{Leveraging Composition and Hierarchy:} The principle of building highly complex functions by stacking or composing simpler computational units. This hierarchical approach allows models to learn features at multiple levels of abstraction, a cornerstone of deep learning.
  \item \textbf{Scaling with Function Approximation:} The strategy of replacing discrete, tabular representations of knowledge (which suffer from the curse of dimensionality) with continuous, parameterized function approximators, such as neural networks, that can generalize across vast state or feature spaces.
\end{itemize}

This report will navigate eight major conceptual pathways, each illustrating one or more of these evolutionary themes. We will begin with the foundational bridge from linear regression to deep neural networks, then explore the progression of unsupervised representation learning from Principal Component Analysis (PCA) to Variational Autoencoders (VAEs). Subsequent sections will chart the course from deterministic to probabilistic clustering, from linear to learned interactions in recommender systems, and through the architectural evolution of memory in sequence models. We will then examine the scaling of reinforcement learning from simple choices to complex policies, the generalization of Bayesian models from naive assumptions to rich dependency structures, and finally, the unification of inference techniques from Expectation-Maximization (EM) to the broader framework of Variational Inference (VI). By tracing these connections, a coherent and unified picture of machine learning emerges, transforming a catalog of algorithms into a story of scientific and engineering progress.

\section{From Linear Regression to Deep Neural Networks: The Dawn of Non-Linearity}

The most fundamental conceptual bridge in machine learning is the one connecting classical linear models to modern deep neural networks. This pathway demonstrates how the core principles of composition and non-linearity transform a simple, interpretable statistical tool into a powerful, universal function approximator. A neural network is not an entirely alien concept but rather a direct and logical extension of the Generalized Linear Model (GLM) framework familiar to statisticians.\refnum{1}

\subsection{The Single Neuron as a Generalized Linear Model (GLM)}
The basic computational unit of a neural network, the neuron, is mathematically equivalent to a GLM. A neuron performs two operations: first, it computes a linear combination of its inputs $\mathbf{x}$ with weights $\mathbf{w}$, plus a bias term $b$, resulting in a linear predictor $z=\mathbf{w}^\top\mathbf{x}+b$.\refnum{2} This is the mathematical core shared by all linear models.

The nature of the activation function determines the specific statistical model to which the neuron is equivalent.
\begin{itemize}[leftmargin=1.2em]
\item \textbf{Linear Regression Equivalence:} If the neuron uses the identity activation, $f(z)=z$, the computation becomes $\hat{y}=\mathbf{w}^\top\mathbf{x}+b$, identical to multivariate linear regression.\refnum{1,4}
\item \textbf{Logistic Regression Equivalence:} If the neuron uses the sigmoid $\sigma(z)=1/(1+e^{-z})$, it becomes equivalent to logistic regression: the sigmoid inverts the logit link, mapping $( -\infty,\infty)$ to $(0,1)$.\refnum{1}
\end{itemize}
This direct correspondence demystifies the fundamental building block of neural networks and anchors them to classical statistics.\refnum{1}

\subsection{The Futility of Linearity: Why Stacking Linear Layers Fails}
If activations remain linear, depth adds no expressivity. For two linear layers $W_1,W_2$, the composition yields $\hat{y}=W_2(W_1\mathbf{x})=(W_2W_1)\mathbf{x}$, which collapses to a single linear layer.\refnum{1,4} Thus depth \emph{requires} nonlinearity.

\subsection{Unlocking Complexity: The Role of Non-Linear Activation Functions}
With nonlinearities (ReLU, Tanh, Sigmoid), the network computes nested nonlinear maps like $f_2(W_2\,f_1(W_1\mathbf{x}+b_1)+b_2)$, which cannot be reduced to a single affine map.\refnum{1} The Universal Approximation Theorem guarantees a single hidden layer with a nonlinear activation can approximate any continuous function on compact subsets of $\R^n$ to arbitrary accuracy.\refnum{1} Practically, deeper networks often learn complex functions more efficiently than shallow wide ones.\refnum{5}

Historically, sigmoids suffered from vanishing gradients; ReLU, $f(z)=\max(0,z)$, alleviates this by maintaining non-saturating derivatives for $z>0$, enabling deep training.\refnum{7,8}

\section{From Data Projection to Generative Manifolds: The Evolution of Unsupervised Representation}

\subsection{The Foundation: Principal Component Analysis (PCA)}
PCA finds an orthogonal basis maximizing projected variance via the eigendecomposition of the covariance matrix; eigenvectors are principal components, eigenvalues quantify explained variance.\refnum{9,11,12}

\subsection{The Bridge: PCA as a Linear Autoencoder}
A one-hidden-layer linear autoencoder with MSE learns the same \emph{principal subspace} as PCA (for the bottleneck dimension $k$), though not necessarily an orthonormal basis.\refnum{14,16} PCA solves analytically; the linear AE learns via gradient descent.\refnum{16}

\subsection{Adding Power: Deep and Non-Linear Autoencoders}
Stacking layers and using nonlinear activations yields deep AEs that learn nonlinear manifolds (e.g., Swiss roll) beyond linear subspaces.\refnum{9,19,8}

\subsection{The Probabilistic Turn: From PCA to PPCA}
PPCA posits $p(\mathbf{z})=\mathcal{N}(0,I)$, $p(\mathbf{x}\mid\mathbf{z})=\mathcal{N}(W\mathbf{z}+\boldsymbol{\mu},\sigma^2 I)$, giving a generative model that handles likelihoods, missing data, and latent structure.\refnum{22,23}

\subsection{The Grand Unification: PPCA as a Linear VAE}
Constrain a VAE to a linear decoder with Gaussian likelihood and standard normal prior, and it reduces to PPCA.\refnum{24} Relaxing linearity to deep nets yields modern VAEs trained by maximizing the ELBO:
\[
\log p_\theta(\mathbf{x}) \ge
\mathbb{E}_{q_\phi(\mathbf{z}\mid\mathbf{x})}\!\big[\log p_\theta(\mathbf{x}\mid\mathbf{z})\big]
- \mathrm{KL}\!\big(q_\phi(\mathbf{z}\mid\mathbf{x})\,\|\,p(\mathbf{z})\big),
\]
optimized via the reparameterization trick $\mathbf{z}=\boldsymbol{\mu}_\phi(\mathbf{x})+\boldsymbol{\sigma}_\phi(\mathbf{x})\odot\boldsymbol{\epsilon}$, $\boldsymbol{\epsilon}\sim\mathcal{N}(0,I)$.\refnum{27}
Posterior collapse links to spurious optima even in linear VAEs/PPCA, motivating objective variants (e.g., $\beta$-VAE).\refnum{24,25,31,32}

\begin{table}[h!]
\centering
\caption{Summary of the PCA $\to$ VAE pathway}
\begin{tabular}{@{}lllll@{}}
\toprule
\textbf{Model} & \textbf{Linearity} & \textbf{Probabilistic} & \textbf{Methodology} & \textbf{Representation} \\
\midrule
PCA & Linear & No  & Eigendecomposition & Orthogonal principal subspace \\
Linear AE & Linear & No & Gradient descent & Basis of principal subspace \\
Deep AE & Nonlinear & No & Gradient descent & Nonlinear manifold \\
PPCA & Linear & Yes & EM / closed-form & Gaussian latent space \\
VAE & Nonlinear & Yes & Variational inference & Flexible latent manifold \\
\bottomrule
\end{tabular}
\end{table}

\section{From Hard Partitions to Probabilistic Mixtures: The Nuances of Clustering}

\subsection{K-Means Clustering: The Deterministic Approach}
K-means partitions data into $K$ clusters minimizing within-cluster sum of squares, yielding hard assignments; assumptions implicit: spherical, similarly sized clusters.\refnum{33,35,36,37}

\subsection{Gaussian Mixture Models: The Probabilistic Generalization}
GMMs model $p(\mathbf{x})=\sum_{k=1}^K \pi_k \,\mathcal{N}(\mathbf{x}\mid\boldsymbol{\mu}_k,\Sigma_k)$, providing soft responsibilities and flexible shapes via full covariances.\refnum{38,35,37,34} Parameters are fit by EM.\refnum{36,42}

\subsection{The Algorithmic Bridge: K-Means as a Limit of EM for GMMs}
With shared spherical covariance $\Sigma_k=\sigma^2 I$ and $\sigma^2\!\to\!0$, EM responsibilities harden to 0/1, and EM steps reduce to k-means assignment/centroid updates.\refnum{33}

\section{From Linear Interactions to Learned Similarities: Recommender Systems}

\subsection{Matrix Factorization (MF)}
Decompose $R\approx P Q^\top$; predict $\hat{r}_{ui}=\mathbf{p}_u^\top\mathbf{q}_i$.\refnum{44,46,49}

\subsection{Neural Collaborative Filtering (NCF)}
Replace the dot product with an MLP: $\hat{r}_{ui}=f(\mathbf{p}_u,\mathbf{q}_i\mid\Theta)$; GMF and NeuMF show MF is a special case and hybrids can capture linear+nonlinear interactions.\refnum{50,49,44} Later work shows tuned MF can rival NCF, encouraging hybrid designs (NeuMF, DeepFM, xDeepFM).\refnum{53,55,56}

\section{The Architecture of Memory: Evolution of Sequence Models}

\subsection{Markov Chains}
First-order Markov property yields one-step memory; limited for complex sequences.\refnum{57}

\subsection{Hidden Markov Models (HMMs)}
Introduce discrete latent states with Markovian transitions and emission distributions; more expressive than observable Markov models.\refnum{60,61,63}

\subsection{Recurrent Neural Networks (RNNs)}
Maintain continuous hidden state $h_t=f(W_{hh}h_{t-1}+W_{xh}x_t)$, a distributed memory; suffer vanishing/exploding gradients.\refnum{64,66,68,7}

\subsection{LSTMs and GRUs}
Gated architectures (input/forget/output gates; cell state $c_t$) preserve gradients, enabling long-range dependencies; GRU simplifies with an update gate.\refnum{7,60,64}

\subsection{Transformers}
Self-attention replaces recurrence, enabling $O(1)$ path length for dependencies at $O(L^2)$ compute, with parallelism and strong long-range modeling.\refnum{66,70,7}

\begin{table}[h!]
\centering
\caption{Key sequence models and memory mechanisms}
\begin{tabular}{@{}lllll@{}}
\toprule
\textbf{Model} & \textbf{Core Concept} & \textbf{Memory} & \textbf{Innovation} & \textbf{Limitation Addressed} \\
\midrule
Markov Chain & Prob. transitions & Memoryless & State dependence & Baseline \\
HMM & Discrete latents & Markov latent & Latent causes & Observable-only limits \\
RNN & Recurrence & Continuous $h_t$ & Distributed memory & Discrete-state limits \\
LSTM/GRU & Gating + $c_t$ & Gated memory & Gradient flow & Long-range credit assignment \\
Transformer & Self-attention & Global context & Parallelism, $O(1)$ path & Sequential bottleneck \\
\bottomrule
\end{tabular}
\end{table}

\section{From Simple Choices to Complex Policies: The Trajectory of RL}

\subsection{Multi-Armed Bandits (MAB)}
Single-state decision-making with exploration--exploitation; a one-state MDP.\refnum{75}

\subsection{Markov Decision Processes (MDPs)}
States $\mathcal{S}$, actions $\mathcal{A}$, transitions $P(s'|s,a)$, rewards $R$, discount $\gamma$; introduce delayed rewards and planning.\refnum{79,81,83}

\subsection{Tabular Q-Learning}
Learn $Q^*(s,a)$ via Bellman updates with a lookup table; works for small discrete spaces.\refnum{84,86}

\subsection{Deep Q-Networks (DQN)}
Use neural nets to approximate $Q(s,a;\theta)$, enabling high-dimensional inputs; stabilize with experience replay and target networks.\refnum{94,85,86}

\section{From Naive Assumptions to Rich Dependencies: Bayesian Models}

\subsection{Naive Bayes}
Assume conditional independence: $p(x_1,\dots,x_n\mid C)=\prod_{i}p(x_i\mid C)$; efficient and surprisingly effective but often violated.\refnum{98,100,99}

\subsection{Bayesian Networks}
DAG + CPTs encode conditional independencies; highly expressive graphical models.\refnum{102,100}

\subsection{The Bridge}
Naive Bayes is a Bayes net with root $C$ and edges $C\to x_i$, no feature--feature edges, making features independent given $C$; TAN relaxes this with limited feature dependencies.\refnum{100,101}

\section{From Maximum Likelihood to Full Distributions: Generalizing Inference}

\subsection{Expectation--Maximization (EM)}
For latent-variable models, EM alternates
E-step $Q(\theta\mid\theta^{\text{old}})=\mathbb{E}_{Z\mid X,\theta^{\text{old}}}[\log p(X,Z\mid\theta)]$ and
M-step $\theta^{\text{new}}=\arg\max_\theta Q(\theta\mid\theta^{\text{old}})$, monotonically increasing $\log p(X\mid\theta)$.\refnum{104}

\subsection{Variational Inference (VI)}
Approximate posteriors by optimizing the ELBO
\[
\log p(X)=\mathcal{L}(q)+\mathrm{KL}\!\big(q(Z)\,\|\,p(Z\mid X)\big),\quad
\mathcal{L}(q)=\mathbb{E}_q[\log p(X,Z)]-\mathbb{E}_q[\log q(Z)],
\]
thus maximizing $\mathcal{L}(q)$ minimizes KL.\refnum{23,106}

\subsection{The Bridge: EM as Special-Case VI}
Coordinate ascent on the ELBO with (i) exact E-step $q(Z)=p(Z\mid X,\theta)$ and (ii) M-step $\arg\max_\theta\mathbb{E}_q[\log p(X,Z\mid\theta)]$ recovers EM exactly. When the E-step is intractable, mean-field VI (``Variational EM'') optimizes within a restricted $q$ family.\refnum{106,107} VAEs are scalable Variational EM with neural parameterizations and a reparameterization trick for backpropagation.\refnum{28,26}

\section*{Conclusion}
The landscape of machine learning, while vast, is underpinned by a coherent set of evolutionary principles. This report has traversed eight distinct but interconnected conceptual pathways, demonstrating that many of the field's most advanced models can be understood as principled generalizations of simpler, foundational concepts. By repeatedly applying themes of introducing non-linearity, embracing probability, leveraging hierarchical composition, and scaling with function approximation, the discipline has systematically expanded its capacity to model the complexity of the real world.

Understanding these conceptual bridges provides practitioners and researchers with a powerful mental model for diagnosing failure modes and charting principled paths for innovation. The frontier of machine learning is built on continual, creative recombination of a core set of ideas; appreciating these connections remains crucial for making sense of ever-growing complexity.

\bigskip
\noindent\textbf{Works Cited}
\begin{enumerate}[leftmargin=1.2em,label={[\arabic*]}]
\item From Linear Models to Neural Networks | Opening the Black Box. Accessed Sept 19, 2025. \url{https://ryansaxe.com/fundamentals/2021/03/03/LinearNN.html}
\item Goings, J. \emph{Neural Networks by Analogy with Linear Regression}. Accessed Sept 19, 2025. \url{http://joshuagoings.com/2020/05/05/neural-network/}
\item Rahman, M.\ A. \emph{Neural Network is Nothing but a Linear Regression}. Accessed Sept 19, 2025. \url{https://medium.com/@asifurrahmanaust/lesson-3-neural-network-is-nothing-but-a-linear-regression-e05a328a0f23}
\item Cross Validated. Difference between Linear Regression and Neural Network. Accessed Sept 19, 2025. \url{https://stats.stackexchange.com/questions/375996/difference-between-linear-regression-and-neural-network}
\item Cross Validated. Are Neural Networks Linear Models? Accessed Sept 19, 2025. \url{https://stats.stackexchange.com/questions/434906/are-neural-networks-linear-models}
\item GeeksforGeeks. Linear Regression vs Neural Networks. Accessed Sept 19, 2025. \url{https://www.geeksforgeeks.org/machine-learning/linear-regression-vs-neural-networks-understanding-key-differences/}
\item SabrePC Blog. RNNs vs LSTM vs Transformers. Accessed Sept 19, 2025. \url{https://www.sabrepc.com/blog/deep-learning-and-ai/rnns-vs-lstm-vs-transformers}
\item Boehmke, B. \emph{Hands-On Machine Learning with R}, Ch.\ 19. \url{https://bradleyboehmke.github.io/HOML/autoencoders.html}
\item Idrees, H. Autoencoders vs PCA. \url{https://medium.com/@hassaanidrees7/autoencoders-vs-pca-dimensionality-reduction-for-complex-data-e07d4612b711}
\item Bhattacharya, A. PCA vs Kernel PCA. \url{https://medium.com/@abhishek8694/linear-vs-non-linear-dimensionality-reduction-pca-and-kernel-pca-10490f345ba9}
\item Texas A\&M notes: PCA \& Autoencoders. \url{https://people.tamu.edu/~sji/classes/PCA.pdf}
\item GeeksforGeeks. How is Autoencoder different from PCA? \url{https://www.geeksforgeeks.org/machine-learning/how-is-autoencoder-different-from-pca/}
\item Wikipedia. Nonlinear Dimensionality Reduction. \url{https://en.wikipedia.org/wiki/Nonlinear_dimensionality_reduction}
\item Cross Validated. Differences between PCA and Autoencoder. \url{https://stats.stackexchange.com/questions/120080/whatre-the-differences-between-pca-and-autoencoder}
\item Plaut, A. \emph{From Principal Subspaces to Principal Components with Linear Autoencoders}. \url{https://arxiv.org/pdf/1804.10253}
\item Reddit r/MachineLearning. PCA vs AutoEncoders. \url{https://www.reddit.com/r/MachineLearning/comments/1gtng8q/d_pca_vs_autoencoders_for_dimensionality_reduction/}
\item Reddit r/MachineLearning. Do people still use autoencoders for dimension reduction? \url{https://www.reddit.com/r/MachineLearning/comments/hyjlot/research_do_people_still_use_autoencoders_for/}
\item \emph{Frontiers}. Linear and Non-linear Dimensionality Reduction on Hand Kinematics. \url{https://www.frontiersin.org/journals/bioengineering-and-biotechnology/articles/10.3389/fbioe.2020.00429/full}
\item Cross Validated. When to use linear vs nonlinear DR? \url{https://stats.stackexchange.com/questions/304262/how-to-know-when-to-use-linear-dimensionality-reduction-vs-non-linear-dimensiona}
\item Reddit r/statistics. Linear vs Nonlinear DR (Q). \url{https://www.reddit.com/r/statistics/comments/156neaf/linear_vs_nonlinear_dimensionality_reduction_q/}
\item Muaz, U. Autoencoders vs PCA: when to use which? \url{https://medium.com/data-science/autoencoders-vs-pca-when-to-use-which-73de063f5d7}
\item Cross Validated. Regular vs Probabilistic PCA. \url{https://stats.stackexchange.com/questions/240624/what-is-the-difference-between-regular-pca-and-probabilistic-pca}
\item Minka, T.\ P.\ et al. \emph{Factor Analysis, PPCA, VI, and VAE}. \url{https://arxiv.org/pdf/2101.00734}
\item TowardsAI. PPCA and VAEs. \url{https://pub.towardsai.net/a-close-look-at-image-generation-using-p-pca-and-variational-autoencoders-90772560ef89}
\item OpenReview. Posterior Collapse in GLVMs. \url{https://openreview.net/pdf/729562a11b8fe6b0af7244d73dea98ec6c5f8376.pdf}
\item Tomczak, J.\ M. VAE notes. \url{https://jmtomczak.github.io/blog/4/4_VAE.html}
\item Wikipedia. Variational Autoencoder. \url{https://en.wikipedia.org/wiki/Variational_autoencoder}
\item Callh, S. VAEs as Probabilistic Inference. \url{https://sebastiancallh.github.io/post/vae-anatomy/}
\item Reddit r/MachineLearning. Probabilistic PCA explanations. \url{https://www.reddit.com/r/MachineLearning/comments/eoal47/d_probabilistic_pca_and_its_explanation/}
\item Cross Validated. What are VAEs used for? \url{https://stats.stackexchange.com/questions/321841/what-are-variational-autoencoders-and-to-what-learning-tasks-are-they-used}
\item NeurIPS. Posterior Collapse and Non-identifiability. \url{https://proceedings.neurips.cc/paper/2021/file/2b6921f2c64dee16ba21ebf17f3c2c92-Paper.pdf}
\item \emph{Don't Blame the ELBO! A Linear VAE Perspective on Posterior Collapse}. \url{https://arxiv.org/abs/1911.02469}
\item \emph{k-means as a variational EM approximation of GMMs}. \url{https://arxiv.org/pdf/1704.04812}
\item Medium. GMM vs K-Means. \url{https://medium.com/@sabarishds03/a-deep-dive-into-gaussian-mixture-model-vs-k-means-algorithm-for-clustering-6599f9a159ae}
\item Cross Validated. GMM vs K-Means for spherical clusters. \url{https://stats.stackexchange.com/questions/489459/in-cluster-analysis-how-does-gaussian-mixture-model-differ-from-k-means-when-we}
\item Medium. K-Means vs Gaussian Mixture. \url{https://medium.com/@amit25173/k-means-clustering-vs-gaussian-mixture-bec129fbe844}
\item VanderPlas, J. \emph{Python DS Handbook}: GMMs. \url{https://jakevdp.github.io/PythonDataScienceHandbook/05.12-gaussian-mixtures.html}
\item EM for GMM (Telecom Paris). \url{https://perso.telecom-paristech.fr/bonald/documents/gmm.pdf}
\item Duke Notes (Cynthia Rudin): GMM/EM. \url{https://users.cs.duke.edu/~cynthia/CourseNotes/GMMEMNotes.pdf}
\item Scikit-learn: Gaussian Mixture. \url{https://scikit-learn.org/stable/modules/mixture.html}
\item Wikipedia. EM algorithm and GMM. \url{https://en.wikipedia.org/wiki/EM_algorithm_and_GMM_model}
\item IITM. GMM via EM. \url{https://www.cse.iitm.ac.in/~vplab/courses/DVP/PDF/gmm.pdf}
\item Stephens, M. Intro to EM (Five Minute Stats). \url{https://stephens999.github.io/fiveMinuteStats/intro_to_em.html}
\item Medium. MF for Recommenders. \url{https://medium.com/thedeephub/recommendation-systems-with-matrix-factorization-544ef5eae7d3}
\item TDS. MF (SVD) Explained. \url{https://towardsdatascience.com/recommendation-system-matrix-factorization-svd-explained-c9a50d93e488/}
\item PMC. Neural MF++ Recommender. \url{https://pmc.ncbi.nlm.nih.gov/articles/PMC10973760/}
\item Cross Validated. Non-negativity in CF. \url{https://stats.stackexchange.com/questions/234027/why-is-non-negativity-important-for-collaborative-filtering-recommender-systems}
\item Wikipedia. Matrix factorization (recommenders). \url{https://en.wikipedia.org/wiki/Matrix_factorization_(recommender_systems)}
\item Bellogin, A. Reenvisioning NCF vs MF. \url{https://abellogin.github.io/2021/recsys.pdf}
\item Nie, L. \emph{Neural Collaborative Filtering}. \url{https://liqiangnie.github.io/paper/p173-he.pdf}
\item He, X. et al. \emph{Neural Collaborative Filtering}. \url{https://arxiv.org/abs/1708.05031}
\item ResearchGate. NCF (request). \url{https://www.researchgate.net/publication/315876128_Neural_Collaborative_Filtering}
\item \emph{NCF vs MF Revisited}. \url{https://arxiv.org/abs/2005.09683}
\item \emph{NCF vs MF Revisited} (PDF). \url{https://arxiv.org/pdf/2005.09683}
\item ResearchGate. NCF vs MF Revisited. \url{https://www.researchgate.net/publication/346821975_Neural_Collaborative_Filtering_vs_Matrix_Factorization_Revisited}
\item MDPI. Deep Matrix Factorization. \url{https://www.mdpi.com/2076-3417/10/14/4926}
\item Reddit. Is deep learning a Markov chain? \url{https://www.reddit.com/r/MachineLearning/comments/47j8j6/is_deep_learning_a_markov_chain_in_disguise/}
\item StackOverflow. Why use RNNs instead of Markov models? \url{https://stackoverflow.com/questions/45341769/why-should-we-use-rnns-instead-of-markov-models}
\item TutorialsPoint. Markov Chains vs HMMs. \url{https://www.tutorialspoint.com/difference-between-markov-chains-and-hidden-markov-models}
\item arXiv. Comparing LSTM and HMM Language Models. \url{https://arxiv.org/pdf/1907.04670}
\item GeeksforGeeks. Markov Chains vs HMMs. \url{https://www.geeksforgeeks.org/artificial-intelligence/what-is-the-difference-between-markov-chains-and-hidden-markov-models/}
\item StackOverflow. Difference between MC and HMM. \url{https://stackoverflow.com/questions/10748426/what-is-the-difference-between-markov-chains-and-hidden-markov-model}
\item Purdue notes. Markov chains and HMMs. \url{https://www.stat.purdue.edu/~junxie/topic3.pdf}
\item \emph{The Mathematical Engineering of DL} (2021), Ch. 8. \url{https://deeplearningmath.org/sequence-models}
\item Viso.ai. Exploring Sequence Models. \url{https://viso.ai/deep-learning/sequential-models/}
\item Medium. Evolution from RNNs to Transformers. \url{https://medium.com/@sajilkumar7276/the-evolution-of-neural-networks-from-rnns-to-transformers-042a20fb6799}
\item Medium. Beginner's Guide to RNN \& LSTMs. \url{https://medium.com/@humble_bee/rnn-recurrent-neural-networks-lstm-842ba7205bbf}
\item HN. Comparing RNN with Markov chain. \url{https://news.ycombinator.com/item?id=11187545}
\item ResearchGate. RNN vs HMM modeling power. \url{https://www.researchgate.net/publication/339328016_Comparing_the_Modeling_Powers_of_RNN_and_HMM}
\item GeeksforGeeks. RNN vs LSTM vs GRU vs Transformers. \url{https://www.geeksforgeeks.org/deep-learning/rnn-vs-lstm-vs-transformers/}
\item Telnyx. Understanding sequence modeling. \url{https://telnyx.com/learn-ai/sequence-modeling}
\item Wikipedia. Transformer (deep learning). \url{https://en.wikipedia.org/wiki/Transformer_(deep_learning_architecture)}
\item Kaggle discussion. Evolution of Transformers and LLMs. \url{https://www.kaggle.com/discussions/general/458973}
\item Baeldung. From RNNs to Transformers. \url{https://www.baeldung.com/cs/rnns-transformers-nlp}
\item Wikipedia. Multi-armed bandit. \url{https://en.wikipedia.org/wiki/Multi-armed_bandit}
\item Medium. Understanding the MAB problem. \url{https://medium.com/@george.felobes/understanding-the-multi-armed-bandit-problem-a-key-concept-in-reinforcement-learning-cb430094274d}
\item Analytics Vidhya. MAB \& Python. \url{https://www.analyticsvidhya.com/blog/2018/09/reinforcement-multi-armed-bandit-scratch-python/}
\item StackOverflow. Why MAB is one-state MDP. \url{https://stackoverflow.com/questions/60164460/why-the-bandit-problem-is-also-called-a-one-step-state-mdp-in-reinforcement-lear}
\item Medium. MDP basics. \url{https://medium.com/@ngao7/markov-decision-process-basics-3da5144d3348}
\item BuiltIn. Understanding MDP. \url{https://builtin.com/machine-learning/markov-decision-process}
\item TDS. Why MDPs matter in RL. \url{https://towardsdatascience.com/why-does-malkov-decision-process-matter-in-reinforcement-learning-b111b46b41bd/}
\item Wikipedia. Markov decision process. \url{https://en.wikipedia.org/wiki/Markov_decision_process}
\item Brown CS Notes. MDP. \url{https://cs.brown.edu/people/rpatel59/latex-notes/mdp-notes.pdf}
\item Wikipedia. Q-learning. \url{https://en.wikipedia.org/wiki/Q-learning}
\item Quora. Q-learning vs Deep Q-learning vs DQN. \url{https://www.quora.com/What-is-the-difference-between-Q-learning-deep-Q-learning-and-deep-Q-network}
\item Baeldung. Q-learning vs DQL vs DQN. \url{https://www.baeldung.com/cs/q-learning-vs-deep-q-learning-vs-deep-q-network}
\item MLQ.ai. Deep Q-Learning guide. \url{https://blog.mlq.ai/deep-reinforcement-learning-q-learning/}
\item TDS. minDQN tutorial. \url{https://towardsdatascience.com/deep-q-learning-tutorial-mindqn-2a4c855abffc/}
\item Spotintelligence. DQN explainer. \url{https://spotintelligence.com/2025/05/12/deep-q-learning-reinforcement-learning/}
\item Medium. Q-learning and DQN. \url{https://medium.com/data-science/q-learning-and-deep-q-networks-436380e8396a}
\item Milvus. Tabular vs function approximation. \url{https://milvus.io/ai-quick-reference/what-is-the-difference-between-tabular-and-function-approximation-methods-in-reinforcement-learning}
\item ResearchGate. Tabular vs deep Q-learning figure. \url{https://www.researchgate.net/figure/A-conceptual-comparison-of-tabular-Q-learning-top-vs-deep-Q-learning-bottom-In-the_fig2_389311892}
\item Medium. DQN explainer. \url{https://medium.com/@samina.amin/deep-q-learning-dqn-71c109586bae}
\item Medium. DQN as function approximator. \url{https://medium.com/intro-to-artificial-intelligence/deep-q-network-dqn-applying-neural-network-as-a-functional-approximation-in-q-learning-6ffe3b0a9062}
\item StackOverflow. RL vs DL vs DRL. \url{https://stackoverflow.com/questions/50542818/whats-the-difference-between-reinforcement-learning-deep-learning-and-deep-re}
\item RL notes. Q-function approximation. \url{https://gibberblot.github.io/rl-notes/single-agent/function-approximation.html}
\item Mnih et al. \emph{Human-level control through deep RL}. \url{https://storage.googleapis.com/deepmind-media/dqn/DQNNaturePaper.pdf}
\item IBM. What are Naïve Bayes classifiers? \url{https://www.ibm.com/think/topics/naive-bayes}
\item Khalid, Z. Naïve Bayes \& Bayesian Networks notes. \url{https://www.zubairkhalid.org/ee514/2023/notes06.pdf}
\item StackOverflow. Bayes net vs Naive Bayes. \url{https://stackoverflow.com/questions/12298150/what-is-the-difference-between-a-bayesian-network-and-a-naive-bayes-classifier}
\item Wikipedia. Naive Bayes classifier. \url{https://en.wikipedia.org/wiki/Naive_Bayes_classifier}
\item Quora. Gaussian NB vs Bayesian networks. \url{https://www.quora.com/What-is-the-difference-between-Gaussian-Naive-Bayes-and-Bayesian-networks}
\item Reddit. BN vs Bayes vs Naive Bayes. \url{https://www.reddit.com/r/MachineLearning/comments/41kne0/bayesian_network_vs_bayesian_inference_vs_naives/}
\item Wikipedia. Expectation--maximization algorithm. \url{https://en.wikipedia.org/wiki/Expectation%E2%80%93maximization_algorithm}
\item CMU. Variational Methods (Zoubin). \url{https://www.cs.cmu.edu/~tom/10-702/Zoubin-702.pdf}
\item Javed, S.\ A. VI and EM. \url{https://stillbreeze.github.io/Variational-Inference-and-Expectation-Maximization/}
\item Wikipedia. Variational Bayesian methods. \url{https://en.wikipedia.org/wiki/Variational_Bayesian_methods}
\item Cross Validated. VI vs Variational EM. \url{https://stats.stackexchange.com/questions/581880/what-is-the-difference-between-variational-inference-and-variational-em}
\item Cambridge ML Group. Variational Bayesian EM. \url{https://mlg.eng.cam.ac.uk/zoubin/papers/valencia02.pdf}
\item Reddit. Approachable explanation of Variational Bayes. \url{https://www.reddit.com/r/MachineLearning/comments/10qt8q/looking_for_an_approachable_explanation_of/}
\item \emph{From EM to Autoencoded Variational Bayes}. \url{https://arxiv.org/abs/2010.13551}
\end{enumerate}

\end{document}
